\scanpage{31}
intérieurs correspondent aux formes dans les
cônes qui sont \uline{non dégénérées}, ou encore
(étant toutes \inserted{pré}riemanniennes \unsure{églt|égalmt}) riemanniennes,
et que pour tt $\psi$ de l'intérieur de
$\Phi^+(\varphi)$ le cône $\Phi^+(\psi)$ est égal à $\Phi^+(\varphi)$.

On trouve ainsi dans $\Phi$ un certain
nombre de cônes convexes \inserted{fermés} $P_\alpha$ avec
intérieur, stables par les $u\otimes u$ ($u\in A$)
et dont les intérieurs sont deux à
deux disjoints, \deleted{\unsure{leur}} \unsure{la|leur} réunion \inserted{des intérieurs} étant l'ens
des formes riemanniennes relativement à $\Phi$.

S'il y a un $\varphi$ riemannienne, alors $-\varphi$ est
riemannienne, mais certainement pas dans le
$=$ cône. (les deux cônes sont opposés).

Quel est la signification de l'invariant :
``nb des cônes riemanniens dans $\Phi$''
(qu'on peut aussi considérer comme
un invariant pour une Alg. riemannienne
sur $Q$, ou seulement sur $\mathbb R$ \unsure{---|-,-}].
% archive sheet number: 30
